\documentclass[../../../include/open-logic-section]{subfiles}

\begin{document}

\olfileid{sth}{replacement}{strength}
\olsection{بدل رکھن دے اصول دی طاقت}

اسیں بدل رکھن دے اصول دی طاقت بارے اک سوکھی گل توں شروع کردے آں۔ جے اسیں $\Z$ توں اگے نہ ودھیے، تاں $\omega + \omega$ دے برابر یا اوہتوں وڈے کسے وی فون نوئمان ترتیبی عدد دی موجودگی ثابت نہیں کر سکدے۔

ایہدی وجہ دا اک خاکہ ویکھو۔ $\ZF$ وچ کم کردیاں سیٹ $V_{\omega+\omega}$ نوں ویکھو۔ ایہ سیٹ~$\Z$ دے اک \emph{نمونے} دا دائرہ بن دا اے۔ ایہ ویکھن لئی اسیں فارمولے نوں کسے دائرے \emph{تک محدود کرن} دی کجھ لکھت متعارف کراندے آں۔
\begin{defn}\ollabel{formularelativization}
	کسے وی سیٹ $M$ تے کسے وی فارمولے $ \phi$ لئی، $\phi^M$ اوہ فارمولا اے جیہڑا $\phi$ دے سارے مسوّراں نوں $M$ تک محدود کر کے ملدا اے۔ یعنی «$\lexists[x]$» دی تھاں «$(\lexists[x \in M])$» لکھو، تے
	«$\lforall[x]$» دی تھاں «$(\lforall[x \in M])$» لکھو۔
\end{defn}\noindent
وکھایا جا سکدا اے کہ ہر مسلّمے $\phi$ لئی جیہڑا~$\Z$ دا اے، $ \ZF \vdash
\phi^{V_{\omega+\omega}}$۔ پر \olref[spine][rank]{ordsetrankalpha} مطابق $\omega+\omega$، $V_{\omega+\omega}$ \emph{دا رکن نہیں}۔ سو $\Z$، $\omega+\omega$ دی غیرموجودگی نال سازگار اے۔

ایسے لئی اسیں \olref[ordinals][replacement]{sec} وچ آکھیا سی کہ
\olref[ordinals][ordtype]{thmOrdinalRepresentation} بدل رکھن دے اصول توں بناں ثابت نہیں ہو سکدا۔ کیوں جے~$\Z$ وچ اک صاف لکھی حُسنِ ترتیب تعریف کرنی سوکھی اے جیہدی ترتیب دی قسم، ساڈی سمجھ مطابق، $\omega+\omega$ \emph{ہونی چاہیدی اے}۔ اسیں \olref[ordinals][idea]{sec} وچ ایس دی اک غیر رسمی مثال دتی سی جدوں اسیں قدرتی عدداں اُتے ایہ ترتیب پیش کیتی سی۔
\begin{align*}
	n \lessdot m \text{ اُدوں تے صرف اُدوں جدوں }&\text{یا تاں }n < m\text{ تے }m-n\text{ جفت اے،}\\
	& \text{یا $n$ جفت اے تے $m$ طاق اے۔}
\end{align*}
پر جے $\omega+\omega$ موجود نہیں، تاں ایہ حُسنِ ترتیب کسے وی ترتیبی عدد نال ہم ساخت نہیں۔ سو $\Z$ توں
\olref[ordinals][ordtype]{thmOrdinalRepresentation} ثابت \emph{نہیں ہُندا}۔

ہُن گل الٹ کے ویکھو۔ بدل رکھن دا اصول $\omega+\omega$ دی موجودگی ثابت کرن دیندا اے، تے ایس لئی $V_{\omega+\omega}$ دی موجودگی وی ثابت کرن دیندا اے۔ گل ایتھے ای نہیں مکدی۔ اسیں جیہڑی \emph{وی} حُسنِ ترتیب تعریف کر سکیے، \olref[ordinals][ordtype]{thmOrdinalRepresentation}
دسدا اے کہ اوہدے نال ہم ساخت کوئی $\alpha$ موجود اے، تے ایس لئی $V_\alpha$ وی موجود اے۔ سو ایویں سیدھے ای بدل رکھن دا اصول ایہ پکا کردا اے کہ سیٹاں دا درجہ وار سلسلہ \emph{بہت اُچا} اے۔

اگلے کجھ حصیاں وچ، تے فیر \olref[card-arithmetic][fix]{sec} وچ، سانوں ہور چنگی طرح پتا لگے گا کہ بدل رکھن دا اصول درجہ وار سلسلے نوں \emph{کِنّا} اُچا بناؤن تے مجبور کردا اے۔ ہن لئی سیدھی گل ایہ اے کہ بدل رکھن دے اصول دا جواز \emph{واقعی} دینا پیندا اے!

\end{document}
